%% ============================================================
%%  Section: Data and Descriptive Statistics
%%  Depositor Discipline in Modern Banking (2026-04-03 draft)
%%
%%  NOTE: Move to latex/section_mds/desc-stats/ before compiling.
%%  Tables/figures referenced from ../docs/tables/ and ../docs/figures/.
%% ============================================================

\section{Data and Sample}
\label{sec:data}

\subsection{Sample Construction}
\label{sec:data:sample}

Our primary sample is drawn from quarterly Call Report filings (FFIEC 041/051)
for all federally insured commercial banks and savings institutions from
2016Q1 through 2025Q4. We restrict to community banks, defined as institutions
with total assets below \$10~billion as of the calendar quarter immediately
preceding their CECL adoption window (2022Q4). This threshold follows the
Federal Reserve's standard community bank definition and excludes the large
SEC-filing institutions that were required to adopt CECL in 2020Q1, whose
adoption timing coincides with the COVID-19 pandemic and would confound our
identification.\footnote{We use the 2020 SEC-filer cohort separately to validate
  the information content of the CECL shock; see
  Section~\ref{sec:identification:signal}.} Banks with no detectable CECL
adoption---defined as a first non-missing CECL adjustment in the eligible
six-quarter window from 2022Q3 through 2023Q4---are excluded, as are
institutions that underwent mergers, conversions, or regulatory actions that
would contaminate the panel within the event window. The final estimation sample
contains 4,320 community banks and approximately 97,000 bank-quarter
observations in the $\pm$10-quarter event window around adoption.

\subsection{The CECL Adjustment: Distribution and Adoption Timing}
\label{sec:data:cecl}

The treatment variable is the CECL Day-One retained-earnings adjustment scaled
by pre-adoption book equity (\texttt{cecl\_equity}), measured at each bank's
adoption quarter. Figure~\ref{fig:cecl_equity_distribution} plots its
cross-sectional distribution. The distribution is heavily right-skewed: the
majority of community banks recorded modest or zero adjustments, reflecting
loan portfolios with limited expected credit losses under the forward-looking
CECL methodology relative to legacy incurred-loss reserves. However, a
substantial right tail of institutions required larger charges, with roughly
five percent of the sample recording an adjustment exceeding 2.5 percent of
pre-adoption equity---our binary ``High CECL'' threshold. The mass near zero
reflects banks with well-seasoned, low-risk loan books for which the incurred-loss
and expected-loss reserves converged; the right tail captures banks whose
portfolios carried credit risk that had been systematically underreserved under
the old standard.

Figure~\ref{fig:cecl_adoption_quarter_distribution} shows the distribution of
adoption dates across the six eligible quarters. Adoption is sharply concentrated
in 2023Q1, which accounts for 3,997 of 4,320 sample banks (92.5\%)---exactly as
the regulatory calendar prescribed for non-public calendar-year filers whose
fiscal year begins January~1. The remaining 323 banks adopted in the surrounding
five quarters, reflecting fiscal-year variation, early voluntary adoption, and
filing-date rounding in the Call Report data. This near-uniform adoption calendar
strengthens identification: the timing of the information shock is almost
entirely driven by a regulatory deadline set years in advance, not by
bank-specific considerations. We nonetheless exploit the staggered minority for
robustness checks using stacked difference-in-differences estimators.

\subsection{Characteristics of High- and Low-CECL Banks}
\label{sec:data:balance}

Table~\ref{tab:descriptive_sumstats_cecl} compares the 217 High-CECL banks
(CECL Adj./Equity $\geq 2.5\%$) to the 4,103 Low-CECL banks at the adoption
cross-section. Several patterns are informative for both the economics and the
identification.

\paragraph{Loan portfolio risk.}
High-CECL banks enter the adoption quarter with meaningfully weaker loan
quality. Their NPL ratio---nonaccrual loans plus loans past due 90 or more days
as a share of total loans---averages 0.74 percent versus 0.54 percent for Low-CECL
banks, a difference of 20 basis points significant at the 1 percent level.
Existing allowances are also larger (1.55 percent vs.\ 1.35 percent of loans,
$p < 0.01$), confirming that high-CECL banks had already accumulated more
conservative reserves under the incurred-loss standard---yet still required a
materially larger Day-One charge when forced to estimate expected losses over the
full contractual life of their loans. This pattern is precisely what the CECL
standard was designed to surface: latent credit risk that had been
systematically underrecognized because it had not yet met the ``probable loss''
threshold under the legacy standard.

\paragraph{Capital adequacy.}
High-CECL banks are less well-capitalized before adoption, with an
equity-to-assets ratio of 8.9 percent compared with 10.3 percent for Low-CECL
banks ($-1.46$ percentage points, $p < 0.01$). This is consistent with
higher-risk loan portfolios being associated with lower capital buffers, and
important to control for in the panel regressions since the CECL charge
mechanically reduces regulatory capital. Our DiD specifications include the
lagged equity ratio as a control to ensure the treatment coefficient captures
the informational content of the disclosure rather than its incidental capital
consequence.

\paragraph{Deposit structure.}
Critically for identification, the two groups look nearly identical on deposit
composition before adoption. Uninsured time deposits as a share of assets
average 5.62 percent for High-CECL banks and 5.28 percent for Low-CECL banks,
a difference of 33 basis points that is not statistically distinguishable from
zero. The share of total deposits held in uninsured accounts is also similar
(37.5 percent vs.\ 39.4 percent, marginally significant at 10 percent with a
sign that actually disfavors High-CECL banks). This near-equivalence in
pre-adoption deposit structure is central to the parallel-trends assumption
underlying our difference-in-differences design: if depositors were already
differentially withdrawing from High-CECL banks before the disclosure, we would
expect to see systematic pre-adoption differences in deposit composition. We do
not. Event-study pre-trend plots in Section~\ref{sec:results} confirm this
formally. One notable exception is insured time deposits, where High-CECL banks
hold more at baseline (9.1 percent vs.\ 7.3 percent of assets, $p < 0.01$),
plausibly reflecting these banks' greater reliance on interest-bearing retail
funding. This baseline difference is absorbed by bank fixed effects in the panel
regressions.

\paragraph{Financial performance.}
High-CECL banks already earn somewhat less before adoption: return on assets
averages 0.235 percent per quarter versus 0.284 percent ($-0.049$ percentage
points, $p < 0.01$), and interest expense is modestly higher (0.264 percent
vs.\ 0.228 percent of assets, $p < 0.01$). These differences are small in
magnitude and, again, absorbed by bank fixed effects. They do confirm, however,
that High-CECL banks are not simply outliers on a single dimension but represent
a coherent group of institutions with elevated credit risk, constrained capital,
and slightly compressed margins---the profile of banks for which a forward-looking
loss disclosure would carry the most new information.

\subsection{What Predicts the CECL Adjustment?}
\label{sec:data:predictors}

Table~\ref{tab:cecl_predictors_cross_section} estimates cross-sectional OLS
regressions of \texttt{cecl\_equity} on bank characteristics at the adoption
quarter. The specifications add controls in three steps: column~(1) includes
log assets and the equity ratio; column~(2) adds NPL ratio, CRE share, and
C\&I share; column~(3) adds the allowance ratio and the unrealized
securities loss ratio.

The results confirm that CECL is driven by credit risk fundamentals. The NPL
ratio is positively and significantly associated with the size of the adjustment
in both columns~(2) and~(3) (coefficient $+0.064$, $p < 0.01$): banks with
more problem loans at adoption recognized larger expected losses under CECL's
forward-looking methodology. The allowance ratio enters positively and
significantly in column~(3) ($+0.159$, $p < 0.01$), indicating that banks
that had accumulated larger incurred-loss reserves still required further CECL
charges---consistent with the adjustment capturing credit risk that had been
inadequately recognized even by the more prudent reservers under the old
standard. Log assets enters positively throughout, suggesting that larger
community banks had more complex loan portfolios generating larger adjustments.

Strikingly, neither the CRE loan share, the C\&I share, nor the unrealized
securities loss ratio is significantly associated with the CECL adjustment,
even though unrealized losses attracted significant attention in the 2023 banking
stress episode. This dissociation from securities risk confirms that
\texttt{cecl\_equity} is a loan-credit-risk measure, not a proxy for
interest-rate risk or securities portfolio quality.

The overall $R^2$ of these regressions is low---at most 1.2 percent---which is
a feature, not a bug. A highly predictable CECL adjustment would raise concerns
that depositors had already impounded the information before the disclosure date,
attenuating our estimates toward zero. The low predictability from observables
supports the interpretation that the Day-One charge contained genuinely new
information---information that was also unanticipated by equity and CDS markets,
as the signal validation evidence in Section~\ref{sec:identification:signal}
demonstrates. Deposit-structure variables are intentionally excluded from these
regressions: if the CECL adjustment were itself predicted by the pre-existing
composition of insured versus uninsured deposits, our identification strategy
would be compromised. The near-zero correlation between \texttt{cecl\_equity}
and pre-adoption deposit mix, visible in Table~\ref{tab:descriptive_sumstats_cecl},
provides the necessary assurance that the treatment is orthogonal to the outcome
of interest at baseline.
